\input format.tex

%\usepackage{arydshln}
\usepackage{graphicx}
\graphicspath{{cores/}}

\begin{document}

%%% diseases

%%%
\setlength{\arrayrulewidth}{0.5pt}
\fontsize{9.3pt}{17pt}\selectfont
\color{gray2}

\vspace*{0mm}
\begin{center}
{\bf\sanhao 相关疾病风险分析—2型糖尿病}
\end{center}

\medskip

\noindent
{\bf\xiaosihao Tom先生：}


\bigskip
2型糖尿病是一种代谢性疾病，表现为多饮、多尿、多食、乏力和消瘦。近年来发病率呈上升趋势。其发病与遗传因素、肥胖、高热量饮食习惯、久坐少动的生活方式及年龄等因素相关。科学研究指出，肠道微生态可能通过参与宿主的能量代谢，最终促进2型糖尿病的发生与发展。

基于您肠道菌群分析，您2型糖尿病的危险因子评估值偏高，保护因子评估值偏高，表明在其它因素相同的情况下，您体内的肠道菌群对2型糖尿病患病风险的影响是由您整体的健康状况决定的，当您身体健康状况改变时，肠道菌群也会发生变化，保护因子和危险因子的评估值也会相应改变。如果危险因子高于保护因子，可能会增加患2型糖尿病的风险。建议您持续监测肠道健康，保持良好的生活习惯，低脂、低糖、低热量饮食，少吃含糖高及升糖指数高的食物，适量补充益生菌，积极预防糖尿病。

2型糖尿病除了与肠道菌群相关外，还与遗传因素、生活方式及其他环境因素等相关。因此，本结果仅作为您进行全面健康管理的一项参考指标，不能作为临床诊断。

{
\lantxh
\fontsize{9.3pt}{9.8pt}\selectfont
\bigskip
%\begin{tctabularx}{tabularx={m{3.75cm}m{4.9cm}<{\centering}m{1.22cm}<{\centering}m{1.8cm}<{\centering}m{2cm}<{\centering}}}
\begin{tctabularx}{tabularx={m{4.8cm}m{5.2cm}<{\centering}m{4.55cm}<{\centering}}}
%\cellcolor{kelly} &
%\cellcolor{kelly} &
%\cellcolor{kelly} &
%\cellcolor{kelly} &
%\cellcolor{kelly}
&&
\\[-6pt]
\cellcolor{kelly} \raisebox{2.614pt}{\color{black00}\bf 测试项目} &
\cellcolor{kelly} \raisebox{2.614pt}{\color{black00}\bf 分布} &
%\cellcolor{kelly} \raisebox{2.614pt}{\color{black00}\bf 检测值} &
%\cellcolor{kelly} \raisebox{2.614pt}{\color{black00}\bf 参考范围} &
\cellcolor{kelly} \raisebox{2.614pt}{\color{black00}\bf 检测结果}
%{\color{black70}\bf 测试项目} & {\color{black00}\bf 分布} & {\color{black00}\bf 检测值} &{\color{black00}\bf 参考范围} & {\color{black00}\bf 检测结果}
\end{tctabularx}

\vspace*{-4.25mm}
\fontsize{8.8pt}{11pt}\selectfont
%\begin{longtable}{m{3.75cm}m{5cm}<{\centering}m{0cm}@{}m{1.22cm}<{\centering}m{1.8cm}<{\centering}m{2cm}<{\centering}}
\begin{longtable}{m{4.8cm}m{5.2cm}<{\centering}m{0cm}@{}m{4.61cm}<{\centering}}
\hline
%\parbox[c]{\hsize}{\vskip7pt 保护因子\\调节免疫功能、能量代谢平衡，降低2型糖尿病风险 \vskip7pt} & \parbox[c]{\hsize}{\vskip7pt\centerline{\raisebox{-1.5ex}{\includegraphics[width=5cm,height=0.3cm]{zxerxingtangniaobingbaohuyinzi_dis.png}}}\vskip7pt}  &\hspace*{-2.691926474cm}\raisebox{-0.45ex}{\includegraphics{triangle.pdf}}&\liuhao 52.69 & \liuhao 11.77\textasciitilde 84.75 & 调节免疫功能、能量代谢平衡的能力较强 \\
\parbox[c]{\hsize}{\vskip7pt 保护因子\\调节免疫功能、能量代谢平衡，降低2型糖尿病风险 \vskip7pt} & \parbox[c]{\hsize}{\vskip7pt\centerline{\raisebox{-1.5ex}{\includegraphics[width=5cm,height=0.3cm]{zxerxingtangniaobingbaohuyinzi_dis.png}}}\vskip7pt}  &\hspace*{-2.691926474cm}\raisebox{-0.45ex}{\includegraphics{triangle.pdf}} & \begin{minipage}{4.60cm}\begin{center}{偏高 \\ \bahao 调节免疫功能、能量代谢平衡的能力较强 }\end{center} \end{minipage} \\
\hline
%\parbox[c]{\hsize}{\vskip7pt 危险因子\\增加代谢失衡、胰岛素抵抗、2型糖尿病风险 \vskip7pt} & \parbox[c]{\hsize}{\vskip7pt\centerline{\raisebox{-1.5ex}{\includegraphics[width=5cm,height=0.3cm]{zxerxingtangniaobingweixianyinzi_dis.png}}}\vskip7pt}  &\hspace*{-4.279141616cm}\raisebox{-0.45ex}{\includegraphics{triangle.pdf}}&\liuhao 2.29 & \liuhao 0\textasciitilde 16.79 & 增加代谢失衡、胰岛素抵抗风险 \\
\parbox[c]{\hsize}{\vskip7pt 危险因子\\增加代谢失衡、胰岛素抵抗、2型糖尿病风险 \vskip7pt} & \parbox[c]{\hsize}{\vskip7pt\centerline{\raisebox{-1.5ex}{\includegraphics[width=5cm,height=0.3cm]{zxerxingtangniaobingweixianyinzi_dis.png}}}\vskip7pt}  &\hspace*{-4.279141616cm}\raisebox{-0.45ex}{\includegraphics{triangle.pdf}} & \begin{minipage}{4.60cm}\begin{center}{偏高 \\ \bahao 增加代谢失衡、胰岛素抵抗风险 }\end{center} \end{minipage} \\
\hline
\end{longtable}

\noindent
\tikz\draw[green2,fill=green2](0,0) circle[radius=4pt]; 区域表示正常人群的一般水平 \tikz\draw[darkblue,fill=darkblue](0,0) circle[radius=4pt]; 区域表示仍在正常水平，但较一般水平偏高或偏低 %\tikz\draw[red,fill=red](0,0) circle[radius=4pt]; 而红色表示为有风险

}

\bigskip
{\centering {\includegraphics{jieshi.pdf}}}
\newpage

%%%
\setlength{\arrayrulewidth}{0.5pt}
\fontsize{9.3pt}{17pt}\selectfont
\color{gray2}

\vspace*{0mm}
\begin{center}
{\bf\sanhao 相关疾病风险分析—结直肠癌}
\end{center}

\medskip

\noindent
{\bf\xiaosihao Tom先生：}


\bigskip
结直肠癌是常见的消化系统恶性肿瘤，随着生活方式的改变，该病发病率在世界范围内呈逐年上升趋势。结直肠癌是大肠黏膜在环境及遗传等多种致病因素作用下发生的恶性变。近年来，有研究发现，人体内肠道菌群作为重要的环境因素在结直肠癌的发生和发展过程中发挥着重要作用。

基于您肠道菌群分析，您结直肠癌的危险因子评估值偏高，保护因子评估值偏高，表明在其他因素相同的情况下，您体内的肠道菌群对您患结直肠癌的风险影响是由您整体的健康状况决定的，当您身体健康状况改变时，肠道菌群也会变化，保护因子和危险因子的评估值也会相应改变。如果危险因子高于保护因子，有可能增加您患结直肠癌的风险。建议您持续监测肠道健康，保持良好的生活习惯，低脂、低热量、适量优质蛋白、适量粗粮饮食，避免高脂肪、高蛋白、油炸、腌制、烟熏、烧烤、霉变及辛辣刺激性食物，补充益生菌，积极预防结直肠癌。

结直肠癌除了与肠道菌群相关外，还与遗传因素、生活习惯及其他环境因素等相关。因此，本结果仅作为您进行全面健康管理的一项参考指标，不能作为临床诊断。

{
\lantxh
\fontsize{9.3pt}{9.8pt}\selectfont
\bigskip
%\begin{tctabularx}{tabularx={m{3.75cm}m{4.9cm}<{\centering}m{1.22cm}<{\centering}m{1.8cm}<{\centering}m{2cm}<{\centering}}}
\begin{tctabularx}{tabularx={m{4.8cm}m{5.2cm}<{\centering}m{4.55cm}<{\centering}}}
%\cellcolor{kelly} &
%\cellcolor{kelly} &
%\cellcolor{kelly} &
%\cellcolor{kelly} &
%\cellcolor{kelly}
&&
\\[-6pt]
\cellcolor{kelly} \raisebox{2.614pt}{\color{black00}\bf 测试项目} &
\cellcolor{kelly} \raisebox{2.614pt}{\color{black00}\bf 分布} &
%\cellcolor{kelly} \raisebox{2.614pt}{\color{black00}\bf 检测值} &
%\cellcolor{kelly} \raisebox{2.614pt}{\color{black00}\bf 参考范围} &
\cellcolor{kelly} \raisebox{2.614pt}{\color{black00}\bf 检测结果}
%{\color{black70}\bf 测试项目} & {\color{black00}\bf 分布} & {\color{black00}\bf 检测值} &{\color{black00}\bf 参考范围} & {\color{black00}\bf 检测结果}
\end{tctabularx}

\vspace*{-4.25mm}
\fontsize{8.8pt}{11pt}\selectfont
%\begin{longtable}{m{3.75cm}m{5cm}<{\centering}m{0cm}@{}m{1.22cm}<{\centering}m{1.8cm}<{\centering}m{2cm}<{\centering}}
\begin{longtable}{m{4.8cm}m{5.2cm}<{\centering}m{0cm}@{}m{4.61cm}<{\centering}}
\hline
%\parbox[c]{\hsize}{\vskip7pt 保护因子\\参与抑制过敏、控制血清胆固醇水平、调节免疫功能 \vskip7pt} & \parbox[c]{\hsize}{\vskip7pt\centerline{\raisebox{-1.5ex}{\includegraphics[width=5cm,height=0.3cm]{zxjiezhichangaibaohuyinzi_dis.png}}}\vskip7pt}  &\hspace*{-1.74718661cm}\raisebox{-0.45ex}{\includegraphics{triangle.pdf}}&\liuhao 17.51 & \liuhao 0\textasciitilde 26.81 & 增强抑制过敏、调节免疫的能力 \\
\parbox[c]{\hsize}{\vskip7pt 保护因子\\参与抑制过敏、控制血清胆固醇水平、调节免疫功能 \vskip7pt} & \parbox[c]{\hsize}{\vskip7pt\centerline{\raisebox{-1.5ex}{\includegraphics[width=5cm,height=0.3cm]{zxjiezhichangaibaohuyinzi_dis.png}}}\vskip7pt}  &\hspace*{-1.74718661cm}\raisebox{-0.45ex}{\includegraphics{triangle.pdf}} & \begin{minipage}{4.60cm}\begin{center}{偏高 \\ \bahao 增强抑制过敏、调节免疫的能力 }\end{center} \end{minipage} \\
\hline
%\parbox[c]{\hsize}{\vskip7pt 危险因子\\与受体结合、刺激炎症因子分泌、引起炎症反应 \vskip7pt} & \parbox[c]{\hsize}{\vskip7pt\centerline{\raisebox{-1.5ex}{\includegraphics[width=5cm,height=0.3cm]{zxjiezhichangaiweixianyinzi_dis.png}}}\vskip7pt}  &\hspace*{-4.452454082cm}\raisebox{-0.45ex}{\includegraphics{triangle.pdf}}&\liuhao 0.01 & \liuhao 0\textasciitilde 24.76 & 增加炎症因子分泌及炎症反应风险 \\
\parbox[c]{\hsize}{\vskip7pt 危险因子\\与受体结合、刺激炎症因子分泌、引起炎症反应 \vskip7pt} & \parbox[c]{\hsize}{\vskip7pt\centerline{\raisebox{-1.5ex}{\includegraphics[width=5cm,height=0.3cm]{zxjiezhichangaiweixianyinzi_dis.png}}}\vskip7pt}  &\hspace*{-4.452454082cm}\raisebox{-0.45ex}{\includegraphics{triangle.pdf}} & \begin{minipage}{4.60cm}\begin{center}{偏高 \\ \bahao 增加炎症因子分泌及炎症反应风险 }\end{center} \end{minipage} \\
\hline
\end{longtable}

\noindent
\tikz\draw[green2,fill=green2](0,0) circle[radius=4pt]; 区域表示正常人群的一般水平 \tikz\draw[darkblue,fill=darkblue](0,0) circle[radius=4pt]; 区域表示仍在正常水平，但较一般水平偏高或偏低 %\tikz\draw[red,fill=red](0,0) circle[radius=4pt]; 而红色表示为有风险

}

\bigskip
{\centering {\includegraphics{jieshi.pdf}}}
\newpage

%%%
\setlength{\arrayrulewidth}{0.5pt}
\fontsize{9.3pt}{17pt}\selectfont
\color{gray2}

\vspace*{0mm}
\begin{center}
{\bf\sanhao 相关疾病风险分析—炎症性肠病}
\end{center}

\medskip

\noindent
{\bf\xiaosihao Tom先生：}


\bigskip
炎症性肠病以持续性肠道非特异性炎症为特征，表现为反复的腹痛、腹泻、粘液血便，甚至出现各种肠道外表现，如视物模糊、关节疼痛、皮疹等。肠道菌群及其代谢产物、宿主基因易感性及肠道黏膜免疫、精神创伤、环境因素以及生活习惯等多种因素共同参与了炎症性肠病的发病机制。患该病人群的肠道菌群多样性较低，病原菌比例较高，产生有害毒素，引起炎症。

基于您肠道菌群分析，您炎症性肠病的危险因子评估值偏高，保护因子评估值偏高，表明在其他因素相同的情况下，您体内的肠道菌群对您患炎症性肠病的风险影响是由您整体的健康状况决定的，当您身体健康状态改变时，肠道菌群也会变化，保护因子和危险因子的评估值也会相应改变。如果危险因子高于保护因子，有可能增加您患肠炎的风险。建议您持续监测肠道健康，保持良好的生活习惯，低盐、低脂、低糖饮食，少吃生冷、油炸、烧烤及辛辣刺激性食物，补充益生菌，积极预防炎症性肠病。

炎症性肠病除了与肠道菌群相关外，还与遗传因素、生活习惯及其他环境因素等相关。因此，本结果仅作为您进行全面健康管理的一项参考指标，不能作为临床诊断。

{
\lantxh
\fontsize{9.3pt}{9.8pt}\selectfont
\bigskip
%\begin{tctabularx}{tabularx={m{3.75cm}m{4.9cm}<{\centering}m{1.22cm}<{\centering}m{1.8cm}<{\centering}m{2cm}<{\centering}}}
\begin{tctabularx}{tabularx={m{4.8cm}m{5.2cm}<{\centering}m{4.55cm}<{\centering}}}
%\cellcolor{kelly} &
%\cellcolor{kelly} &
%\cellcolor{kelly} &
%\cellcolor{kelly} &
%\cellcolor{kelly}
&&
\\[-6pt]
\cellcolor{kelly} \raisebox{2.614pt}{\color{black00}\bf 测试项目} &
\cellcolor{kelly} \raisebox{2.614pt}{\color{black00}\bf 分布} &
%\cellcolor{kelly} \raisebox{2.614pt}{\color{black00}\bf 检测值} &
%\cellcolor{kelly} \raisebox{2.614pt}{\color{black00}\bf 参考范围} &
\cellcolor{kelly} \raisebox{2.614pt}{\color{black00}\bf 检测结果}
%{\color{black70}\bf 测试项目} & {\color{black00}\bf 分布} & {\color{black00}\bf 检测值} &{\color{black00}\bf 参考范围} & {\color{black00}\bf 检测结果}
\end{tctabularx}

\vspace*{-4.25mm}
\fontsize{8.8pt}{11pt}\selectfont
%\begin{longtable}{m{3.75cm}m{5cm}<{\centering}m{0cm}@{}m{1.22cm}<{\centering}m{1.8cm}<{\centering}m{2cm}<{\centering}}
\begin{longtable}{m{4.8cm}m{5.2cm}<{\centering}m{0cm}@{}m{4.61cm}<{\centering}}
\hline
%\parbox[c]{\hsize}{\vskip7pt 保护因子\\阻止细菌黏附和易位，改善肠黏膜通透性 \vskip7pt} & \parbox[c]{\hsize}{\vskip7pt\centerline{\raisebox{-1.5ex}{\includegraphics[width=5cm,height=0.3cm]{zxchangyanbaohuyinzi_dis.png}}}\vskip7pt}  &\hspace*{-1.137071372cm}\raisebox{-0.45ex}{\includegraphics{triangle.pdf}}&\liuhao 16.67 & \liuhao 0\textasciitilde 21.13 & 阻止细菌黏附和易位，改善肠黏膜通透性能力较高 \\
\parbox[c]{\hsize}{\vskip7pt 保护因子\\阻止细菌黏附和易位，改善肠黏膜通透性 \vskip7pt} & \parbox[c]{\hsize}{\vskip7pt\centerline{\raisebox{-1.5ex}{\includegraphics[width=5cm,height=0.3cm]{zxchangyanbaohuyinzi_dis.png}}}\vskip7pt}  &\hspace*{-1.137071372cm}\raisebox{-0.45ex}{\includegraphics{triangle.pdf}} & \begin{minipage}{4.60cm}\begin{center}{偏高 \\ \bahao 阻止细菌黏附和易位，改善肠黏膜通透性能力较高 }\end{center} \end{minipage} \\
\hline
%\parbox[c]{\hsize}{\vskip7pt 危险因子\\破坏肠道耐受机制，损伤肠上皮细胞和肠黏膜屏障 \vskip7pt} & \parbox[c]{\hsize}{\vskip7pt\centerline{\raisebox{-1.5ex}{\includegraphics[width=5cm,height=0.3cm]{zxchangyanweixianyinzi_dis.png}}}\vskip7pt}  &\hspace*{-3.853294844cm}\raisebox{-0.45ex}{\includegraphics{triangle.pdf}}&\liuhao 0.06 & \liuhao 0\textasciitilde 16.60 & 促进破坏肠上皮细胞和肠黏膜屏障 \\
\parbox[c]{\hsize}{\vskip7pt 危险因子\\破坏肠道耐受机制，损伤肠上皮细胞和肠黏膜屏障 \vskip7pt} & \parbox[c]{\hsize}{\vskip7pt\centerline{\raisebox{-1.5ex}{\includegraphics[width=5cm,height=0.3cm]{zxchangyanweixianyinzi_dis.png}}}\vskip7pt}  &\hspace*{-3.853294844cm}\raisebox{-0.45ex}{\includegraphics{triangle.pdf}} & \begin{minipage}{4.60cm}\begin{center}{偏高 \\ \bahao 促进破坏肠上皮细胞和肠黏膜屏障 }\end{center} \end{minipage} \\
\hline
\end{longtable}

\noindent
\tikz\draw[green2,fill=green2](0,0) circle[radius=4pt]; 区域表示正常人群的一般水平 \tikz\draw[darkblue,fill=darkblue](0,0) circle[radius=4pt]; 区域表示仍在正常水平，但较一般水平偏高或偏低 %\tikz\draw[red,fill=red](0,0) circle[radius=4pt]; 而红色表示为有风险

}

\bigskip
{\centering {\includegraphics{jieshi.pdf}}}
\newpage

\end{document}
